\documentclass[10pt]{article}
\usepackage[margin=0.9in]{geometry}
\usepackage[T1]{fontenc}
\usepackage{lmodern}
\usepackage{microtype}
\usepackage{booktabs,array,tabularx}
\usepackage{enumitem}
\usepackage{xcolor}
\usepackage[colorlinks=true,linkcolor=black,urlcolor=blue!50!black]{hyperref}
\setlist{itemsep=1pt,topsep=3pt,leftmargin=*}
\setlength{\parskip}{4pt}
\setlength{\parindent}{0pt}
\newcommand{\nctx}{365}
\newcommand{\neone}{2,500}
\newcommand{\exactDSFlash}{98}
\newcommand{\bestmodel}{DeepSeek V4 Flash}
\newcommand{\tiedwithbest}{DeepSeek V4 Pro}
\newcommand{\nsigbelow}{three}
\newcommand{\deltamin}{7}
\newcommand{\deltamax}{17}
\newcommand{\exactmin}{81}
\newcommand{\exactmax}{98}
\newcommand{\recallmin}{0.966}
\newcommand{\recallmax}{0.996}
\newcommand{\otherswrong}{25}
\newcommand{\othersmadeup}{0}
\newcommand{\flashmadeup}{1,032}
\newcommand{\flashworstmade}{991}
\newcommand{\flashworstgold}{100}
\newcommand{\medianmissed}{1}
\newcommand{\ninvalid}{7}
\newcommand{\nrefused}{2}
\newcommand{\failDSPro}{4}
\newcommand{\failDSFlash}{2}
\newcommand{\largeHaiku}{53}
\newcommand{\largeMiMoFlash}{34}
\newcommand{\nenum}{100}
\newcommand{\smallmin}{93}
\newcommand{\orpooledanswer}{0.50}
\newcommand{\orpooledanswerci}{0.36--0.60}
\newcommand{\orpooledtable}{0.39}
\newcommand{\orpooledtableci}{0.22--0.56}
\newcommand{\fmtMiMoFlashJson}{86}
\newcommand{\fmtMiMoFlashNl}{73}
\newcommand{\condsizeone}{13}
\newcommand{\condsizetwo}{6.5}
\newcommand{\condsizethree}{3.5}
\newcommand{\ragrecall}{0.869 for all models}
\newcommand{\nscalar}{35}
\newcommand{\scalarerrors}{11}
\newcommand{\scalarunder}{9}
\newcommand{\rsmallDSPro}{1,050}
\newcommand{\rlargeDSPro}{8,832}
\newcommand{\rsmallHaiku}{173}
\newcommand{\rlargeHaiku}{6,442}
\newcommand{\flashshallowexact}{62}
\newcommand{\flashdeepexact}{72}
\newcommand{\etwonote}{DeepSeek V4 Pro's E2 runs are pending (provider usage limits).}
\newcommand{\agreemin}{87}
\newcommand{\agreemax}{100}
\newcommand{\flipmin}{0}
\newcommand{\flipmax}{13}
\newcommand{\reasonshare}{77--97}
\newcommand{\totalcalls}{3,796}
\newcommand{\failedcalls}{1,163}

\pagestyle{empty}

\begin{document}
{\Large\textbf{Memo: Completeness of relational operations in LLMs}}\par
\vspace{2pt}
\textbf{From:} Daryl Alfaro \quad \textbf{Date:} \today \quad
\textbf{Re:} preliminary exploration (deliverables from the research brief)\par
\vspace{-2pt}\rule{\linewidth}{0.4pt}

\textbf{Summary.} The phenomenon is real and measurable. On \nctx{} questions whose records were all in the
prompt, five models had mean recall of \recallmin--\recallmax, yet only \exactmin--\exactmax\% of their answers
were exactly the correct set. The errors are omissions, not hallucinations, for four of the five models.
Completeness falls as answer sets and tables grow, and prompting does not fix it. The two DeepSeek models are near
ceiling at the current difficulty, so a full benchmark needs harder tiers. The attached results paper has the
details, and the code, data and every raw response are at
\url{https://github.com/daryl-888/llm-relational-completeness}.

\textbf{Deliverables against the brief}

\begin{tabularx}{\linewidth}{@{}>{\raggedright\arraybackslash}p{3.1cm}>{\raggedright\arraybackslash}X@{}}
\toprule
Brief asked for & Status \\
\midrule
Related work & Done; summary below, and the full review of about 20 papers is attached
(\texttt{related\_work.pdf}). \\
Proof of concept & Done. Reproducible failures with known ground truth. For example, DeepSeek V4 Pro was incomplete
on the same 50-item list in all four runs, and Claude Haiku 5.5 left Maine out of ``U.S.\ states containing `in'\,''
in all four runs (paper, Box~1 and the failure-case table). \\
Preliminary generalization & Done, and broader than asked. Selection, JOIN and aggregation, plus projection, nested
queries and set operations. Every operation except set operations produces incomplete answers for at least three
models. Wrong totals on COUNT/SUM questions are mostly undercounts (\scalarunder{} of \scalarerrors). \\
Benchmark proposal & Built and run rather than only proposed: 500 questions, a task taxonomy of 7 operations
$\times$ 3 settings, SQL-derived ground truth, and the metrics exact-set accuracy, recall/precision and a
missing/near-miss/wrong/made-up/duplicate error taxonomy. \\
Feasibility assessment & Below. \\
\bottomrule
\end{tabularx}

\textbf{Closest prior work.} \emph{HoloBench} (Maekawa et al., ICLR 2025) is closest. It writes out database rows
as long text and evaluates five query types while varying context length and information density, but it scores
with a model-based judge and partial credit, and does not control the number of correct items. \emph{Oolong}
(Bertsch et al., 2025) covers long-context aggregation with scalar answers. \emph{LongProc} (2025) scores
HTML-to-TSV extraction by row-level F1, with no selection conditions. \emph{LOFT} (Lee et al., 2024) finds
SQL-like compositional reasoning hard for long-context models. Multi-answer QA (QAMPARI, RoMQA, QUEST) mixes up
retrieval, knowledge and reading failures. \emph{Galois} (Saeed et al., EDBT 2024) queries a model's memory with
SQL. What is new here:
\begin{enumerate}
\item answer-set size is a controlled variable;
\item ground truth comes from executing SQL, with a set-level error taxonomy and no LLM judge;
\item one framework covers seven operations;
\item it reports the reasoning cost of completeness.
\end{enumerate}
\textbf{Novelty risk: medium.} HoloBench is the baseline to compare against directly.

\textbf{Feasibility verdict: proceed, with changes.}
\begin{itemize}
\item \emph{Robust?} Yes, as a rate. Two independent samples agree on \agreemin--\agreemax\% of verdicts, and model
differences are statistically clear: \nsigbelow{} models trail the best by \deltamin--\deltamax{} points, with intervals excluding zero.
Individual failures sometimes recur and sometimes don't, so results should be reported over many questions and
several samples.
\item \emph{Generalizable?} Across operations, mostly yes. Across access paths, partly: the memory-only slice
shows the same omissions, but the retrieval slice as built measures only the retriever.
\item \emph{Distinct from existing work?} Yes on answer-size control and set-level error types; that distinction
should be shown empirically against HoloBench.
\item \emph{Headroom?} Limited for the strongest models: the DeepSeek models are exact on \exactDSFlash\% of
in-context questions. MiMo-V2.6 Flash, the weakest, is exact on only \largeMiMoFlash\% of questions with 50 or
more correct items.
\end{itemize}

\textbf{Recommended changes for a full benchmark}
\begin{enumerate}
\item \textbf{Harder tiers:} answer sets of 200--1{,}000 items and tables of 5k--20k rows, to restore headroom.
\item \textbf{A fully crossed design} of answer size $\times$ table size $\times$ format $\times$ number of
conditions. These are partly
confounded today (large answers only appear in large tables) and are currently separated by regression.
\item \textbf{Control reasoning:} a reasoning-off condition and fixed reasoning budgets, to test whether the
reasoning--completeness link is causal.
\item \textbf{Redesign the retrieval slice} with queries where the retriever's top results include near-miss
distractors.
\item \textbf{More rigour:} at least 3 samples per question, more questions per cell (the main trend rests on 100
questions per model), a real-data slice, and frontier models.
\end{enumerate}

\textbf{Questions for you.}
\begin{enumerate}
\item Should the next step target a paper (with a direct HoloBench comparison), or first a larger benchmark?
\item Which models should be in scope, and is there a compute or API budget for frontier models?
\end{enumerate}

\end{document}
